% The PGG framework blueprint: modular monodromy MPC + the den Boer instance.
% Module prefixes:  pgg_smc.<mod>  (lib/protocol/groups/security/instances)
%                   pgg_reconstruct.<mod>  (reconstruct/)
% Status: \rocqok = proved/defined in Rocq (green). No \rocqok = an assumed,
% cited classical fact (the six justified axioms), rendered blue-dashed.

\chapter{Introduction}

The PGG framework runs a card-shuffling multiparty protocol whose recovered
secret lives in a starting layout, with the monodromy cut as anonymizing
randomness. Its guarantees are theorems quantified over interface records; an
\emph{instance} is a tuple of records that fills those interfaces. This
blueprint makes that modularity structural. Part~I introduces every interface
and every generic theorem with no instance named. Part~II is pure
instantiation: den Boer as the worked end-to-end hero, the five-card family of
Kim as the biased anonymity instance, and $S_5$ / $S_5\times S_5$ as the large
groups the admissibility gate rejects as algebraic-geometry codes.

The worked MPC (den Boer), its quantitative leakage analysis (Kim), and the
impossibility theorem (the $S_5$ no-go) are entirely axiom-free. Only the
statement that the large groups exist as covers of the right genus is taken on
trust, as six cited classical facts confined to the $S_5$ and $S_5\times S_5$
chapters. In the dependency graph these six are the only blue-dashed nodes;
everything else is green.

\begin{figure}[t]
\centering
\begin{tikzpicture}[
  iface/.style={draw, rounded corners, align=center, inner sep=4pt, font=\small,
                text width=4.0cm, minimum height=0.7cm, fill=black!4},
  plug/.style={draw, rounded corners, align=center, inner sep=3pt, font=\footnotesize,
               text width=3.2cm, minimum height=0.6cm},
  band/.style={draw, rounded corners, align=center, inner sep=4pt, font=\footnotesize,
               text width=6.4cm, fill=black!3},
  verdict/.style={draw, rounded corners, align=center, inner sep=3pt, font=\footnotesize,
                  text width=4.2cm, minimum height=0.6cm},
  every path/.append style={-stealth, thick}, font=\small ]
\node[iface] (mono)  at (0,9.0) {\textbf{MonodromyRepr} / PGGInterface\\ \scriptsize cut group $G$, starts $\pi$};
\node[iface] (ts)    at (0,7.5) {\textbf{ThresholdScheme}\\ \scriptsize share $\cdot$ reconstruct};
\node[iface] (plug)  at (0,6.0) {\textbf{ReconPlug}\\ \scriptsize reconstruction invariance};
\node[iface] (ie)    at (0,4.5) {\textbf{InputEncoding}\\ \scriptsize inputs $\mapsto$ layout};
\node[iface] (prof)  at (0,3.0) {\textbf{MonodromyProfile} $+$ SecurityWitness\\ \scriptsize anonymity $\eps$};
\node[iface] (pismc) at (0,1.5) {\textbf{piSMC dealer} $+$ session types};
\node[plug] (fkim)  at (6.2,9.0) {$\mathtt{FiveCardKim\_M}$, $\pi=\mathrm{ord\_tuple}\,5$};
\node[plug] (fcs)   at (6.2,7.5) {$\mathtt{fcI\_scheme}$};
\node[plug] (fcp)   at (6.2,6.0) {$\mathtt{five\_card\_plug}$};
\node[plug] (dbe)   at (6.2,4.5) {$\mathtt{den\_boer\_encoding}$};
\node[plug] (dbp)   at (6.2,3.0) {$\mathtt{den\_boer\_profile}$ ($\eps=0$)};
\node[plug] (dbd)   at (6.2,1.5) {$\mathtt{den\_boer\_dealer\_layout}$};
\draw (fkim) -- (mono);  \draw (fcs) -- (ts);    \draw (fcp) -- (plug);
\draw (dbe)  -- (ie);    \draw (dbp) -- (prof);  \draw (dbd) -- (pismc);
\node[band] (gen) at (0,-0.2) {\textbf{generic theorems} (proven once):\\
  $\mathtt{pgg\_recon\_monodromy\_correct}$ $\cdot$ $\mathtt{recon\_from\_layout\_output}$\\
  $\cdot$ $\mathtt{ts\_private}$ $\cdot$ $\mathtt{sw\_bound}$ $\cdot$ session duality};
\draw (ts) -- (gen); \draw (plug) -- (gen); \draw (ie) -- (gen);
\draw (prof) -- (gen); \draw (pismc) -- (gen);
\node[band, text width=4.0cm, fill=black!1] (e2e) at (6.2,-0.2)
  {\textbf{den Boer end-to-end:}\\ recovers $a\AND b$, $\MI(\text{in};\text{view}\mid\text{out})=0$,\\
   $\mathtt{leak\_k}$ ramp, duality};
\draw (dbd) -- (e2e); \draw (dbe) -- (e2e);
\node[band, text width=7.2cm] (gate) at (3.0,-2.2)
  {\textbf{AlgebraicRigidity gate} \;($\mathtt{ar\_large\_group\_forces\_gap}\cdot\mathtt{ar\_gap\_bound}$)\\
   Klein cap $\Klein$ sorts instances by $|G|$};
\draw (gen) -- (gate);
\node[verdict, fill=green!8] (sharp) at (-0.3,-4.0)
  {$|G|\le 60$: genus-0 sharp\\ \textbf{den Boer} $C_5$, $|G|=5$ \;$\checkmark$ AG code};
\node[verdict, fill=red!6]  (nogo)  at (6.0,-4.0)
  {$|G|>60$: AG-code no-go\\ \textbf{s5} $|G|{=}120$ $\times$, \textbf{s5x5} $|G|{=}14400$ $\times$\\
   \scriptsize genus forced $>0$; s5 $\to$ Bring genus 4; gap vacuous};
\draw (gate) -- (sharp); \draw (gate) -- (nogo);
\node[font=\itshape\footnotesize, anchor=south] at (0,9.7) {Part I: the framework (no instance named)};
\node[font=\itshape\footnotesize, anchor=south] at (6.2,9.7) {Part II: instantiation};
\end{tikzpicture}
\caption{The PGG framework's modularity. The interface column (left, Part~I)
names no instance; an instance is exactly a tuple of arrows filling the six
slots. den Boer (right, Part~II) fills all six and inherits every generic
theorem (band, centre): reconstruction is monodromy-invariant, the input
encoding's output is correct under every cut, the scheme is threshold-private,
and the dealing is anonymous. Its end-to-end guarantees follow. The
admissibility gate (bottom) sorts monodromy groups by the Klein cap $\Klein$:
small groups admit a sharp genus-0 AG code (den Boer's $C_5$, $\checkmark$), while
large groups are forced off genus 0, where the threshold gap is vacuous, so they
are an AG-code no-go (s5, s5x5, $\times$). The genus and gap are the AG-code
lens; the $\checkmark$/$\times$ is the verdict.}
\label{fig:bp:modularity}
\end{figure}

\part{The modular framework}

\chapter{The monodromy interface}

\begin{definition}[Monodromy representation]
  \label{def:monodromy_repr}
  \rocq{pgg_smc.pgg_interface.MonodromyReprType,
       pgg_smc.pgg_interface.PGGTypes,
       pgg_smc.pgg_interface.MonodromyReprWithGeneratorType}
  \rocqok
  A monodromy representation packages a finite group $G$ acting by a morphism
  $\mono : G \to \mathfrak{S}_N$ on $N$ card positions. The generator-aware
  variant adds a generating tuple, supporting the word search of the security
  analysis.
\end{definition}

\begin{definition}[Protocol interface and starting layout]
  \label{def:pgg_interface}
  \rocq{pgg_smc.pgg_interface.PGGInterface,
       pgg_smc.pgg_interface.pgg_dtype,
       pgg_smc.pgg_interface.pgg_data}
  \rocqok
  \uses{def:monodromy_repr}
  A protocol interface fixes the player count $T$ and a tuple $\pi$ of distinct
  starting positions. The session data type tags the three message kinds dealt
  during the protocol.
\end{definition}

\begin{definition}[Endpoints and the word layer]
  \label{def:endpoint}
  \rocq{pgg_smc.pgg_interface.endpoint,
       pgg_smc.pgg_interface.endpoints,
       pgg_smc.pgg_interface.word_eval,
       pgg_smc.pgg_interface.search_space}
  \rocqok
  \uses{def:monodromy_repr,def:pgg_interface}
  The endpoint of a cut $\cut$ at a position is its image $\mono(\cut)$.
  A length-$L$ word evaluates to a group element, and the search space is the
  cardinality of the reachable set, bounded by $|G|$.
\end{definition}

\chapter{Sharing, reconstruction, and the recon plug}

\begin{definition}[Threshold scheme]
  \label{def:threshold_scheme}
  \rocq{pgg_reconstruct.pgg_sharing_framework.ThresholdScheme}
  \rocqok
  A threshold scheme bundles a validity predicate $\val$, a reconstruction map
  $\recon$, and an encoder, with correctness ($\val~\secret~\text{shares}
  \Rightarrow \recon~\text{shares}=\secret$) and privacy (sub-threshold
  coalitions are indistinguishable across secrets) as axioms.
\end{definition}

\begin{definition}[Reconstruction invariance and the recon plug]
  \label{def:recon_plug}
  \rocq{pgg_reconstruct.pgg_sharing_framework.ts_recon_perm_invariant,
       pgg_reconstruct.covering_scheme.ReconPlug}
  \rocqok
  \uses{def:threshold_scheme,def:monodromy_repr}
  Reconstruction invariance is the position-permutation law
  $\recon[\,\text{tnth shares}~(\sigma(g)~i)\,] = \secret$ for $g\in G$. A recon
  plug bundles a scheme, a content readout, the monodromy-induced permutation,
  and a proof of this invariance over the full group.
\end{definition}

\begin{definition}[Endpoint reconstruction]
  \label{def:pgg_recon}
  \rocq{pgg_reconstruct.pgg_sharing_framework.pgg_recon,
       pgg_reconstruct.pgg_sharing_framework.pgg_recon_endpoints}
  \rocqok
  \uses{def:threshold_scheme,def:pgg_interface}
  Reconstruction of the dealt endpoints reads each player's shuffled start
  through the content map and applies $\recon$.
\end{definition}

\begin{theorem}[Recovery is monodromy-invariant]
  \label{thm:recon_monodromy_correct}
  \rocq{pgg_reconstruct.pgg_sharing_framework.pgg_recon_monodromy_correct}
  \rocqok
  \uses{def:pgg_recon,def:recon_plug}
  For any hidden cut $P$ in the monodromy group, reconstructing the endpoints
  recovers the secret. This is the central correctness guarantee of the whole
  framework: it depends only on reconstruction invariance and a valid starting
  layout, not on any instance.
\end{theorem}

\begin{definition}[Covering data and the covering scheme]
  \label{def:covering_scheme}
  \rocq{pgg_reconstruct.covering_scheme.CoveringData,
       pgg_reconstruct.covering_scheme.CoveringScheme,
       pgg_reconstruct.covering_scheme.genus_from_hurwitz}
  \rocqok
  \uses{def:recon_plug}
  Covering data records the Riemann--Hurwitz parameters of a curve cover and
  determines the covering genus $\genus$. A covering scheme couples a recon plug
  to covering data with the gap bound $T \le k + 2\genus$.
\end{definition}

\chapter{Input encoding}

\begin{definition}[Input encoding]
  \label{def:input_encoding}
  \rocq{pgg_reconstruct.input_encoding.InputEncoding}
  \rocqok
  \uses{def:recon_plug}
  An input encoding is the deterministic half of a randomized encoding of a
  function $f$ over a plug. It maps each input to a valid share layout
  ($\mathtt{ie\_assemble}$), names the computed function $f$
  ($\mathtt{ie\_output}$), and requires equal-output inputs to lie in one
  cut-orbit ($\mathtt{ie\_orbit}$). The cut supplies the randomness.
\end{definition}

\begin{theorem}[Output correctness for every cut]
  \label{thm:ie_output_correct}
  \rocq{pgg_reconstruct.input_encoding.ie_output_correct}
  \rocqok
  \uses{def:input_encoding,def:recon_plug}
  For every cut $g_0$ in the group, reconstructing the cut-permuted assembled
  layout returns $f$ of the input. The output is independent of the cut.
\end{theorem}

\begin{theorem}[Layout recovery]
  \label{thm:recon_from_layout_output}
  \rocq{pgg_reconstruct.input_encoding.recon_from_layout,
       pgg_reconstruct.input_encoding.recon_from_layout_output}
  \rocqok
  \uses{thm:ie_output_correct}
  The reindex-form recovery of an encoded input's layout returns $f$, generic
  over the plug and the encoded function. The den Boer output theorem is this
  result at one instance.
\end{theorem}

\chapter{Anonymity: profile and security witness}

\begin{definition}[Security witness]
  \label{def:security_witness}
  \rocq{pgg_reconstruct.algebraic_rigidity.SecurityWitness,
       pgg_reconstruct.algebraic_rigidity.SecurityExact,
       pgg_reconstruct.algebraic_rigidity.SecurityAsymptotic}
  \rocqok
  A security witness carries an always-present variation-distance bound on the
  cut distribution, optionally an exact equality, and optionally an asymptotic
  spectral-gap convergence certificate. The bound is the anonymity parameter
  $\eps$.
\end{definition}

\begin{definition}[Monodromy profile]
  \label{def:monodromy_profile}
  \rocq{pgg_smc.pgg_monodromy_profile.MonodromyProfile}
  \rocqok
  \uses{def:monodromy_repr,def:pgg_interface,def:security_witness,def:recon_plug}
  A monodromy profile is a fully plugged group program: a monodromy
  representation, a protocol interface, a security witness, and a recon plug.
\end{definition}

\chapter{The piSMC realization}

\begin{definition}[The exchange programs]
  \label{def:exchange_programs}
  \rocq{pgg_smc.card_exchange_pismc.exchange_dealer,
       pgg_smc.card_exchange_pismc.exchange_player,
       pgg_smc.card_exchange_pismc.exchange_verifier,
       pgg_smc.pgg_input_commitment.exchange_dealer_with_commit}
  \rocqok
  \uses{def:pgg_interface}
  The dealer deals each player a hand and a selection index; the player reveals
  one card position; the verifier collects them. A commit prologue prepends one
  receive per input party, threading the committed values into the dealing body.
\end{definition}

\begin{lemma}[Session-type duality]
  \label{lem:session_duality}
  \rocq{pgg_smc.card_exchange_pismc.channels_dual,
       pgg_smc.card_exchange_pismc.dealer_player0_dual_gen,
       pgg_smc.pgg_input_commitment.exchange_dealer_with_commit_nil}
  \rocqok
  \uses{def:exchange_programs}
  Every pair of protocol processes has dual session types, discharged by
  reduction. The check is on the channel structure, independent of the dealt
  payload, so it transfers to any content readout.
\end{lemma}

\chapter{Admissibility: the algebraic-rigidity gate}

\begin{definition}[The Klein genus-0 bound]
  \label{def:klein_bound}
  \rocq{pgg_reconstruct.cover_tradeoff.klein_genus0_bound,
       pgg_reconstruct.cover_tradeoff.genus0_automorphism_bound}
  \rocqok
  \uses{def:covering_scheme}
  A genus-0 cover embeds $G$ in $\mathrm{PGL}(2,\overline{\GF})$, so Klein's
  classification bounds $|G| \le \max(2N,60)$, with $60=|A_5|$ the largest
  polyhedral exception. The constant is the Klein cap, not $|\mathrm{PGL}(2,q)|$.
\end{definition}

\begin{definition}[Algebraic rigidity]
  \label{def:algebraic_rigidity}
  \rocq{pgg_reconstruct.algebraic_rigidity.ThresholdWitness,
       pgg_reconstruct.algebraic_rigidity.AlgebraicRigidity}
  \rocqok
  \uses{def:security_witness,def:covering_scheme,def:klein_bound}
  Algebraic rigidity combines a security witness with a threshold witness, the
  latter carrying the genus-0 obligation $\genus=0 \Rightarrow |G| \le
  \mathtt{klein\_genus0\_bound}$.
\end{definition}

\begin{theorem}[The rigidity laws]
  \label{thm:rigidity_laws}
  \rocq{pgg_reconstruct.algebraic_rigidity.ar_large_group_forces_gap,
       pgg_reconstruct.algebraic_rigidity.ar_gap_bound,
       pgg_reconstruct.algebraic_rigidity.ar_genus_gap_dichotomy,
       pgg_reconstruct.algebraic_rigidity.ar_protocol_correct}
  \rocqok
  \uses{def:algebraic_rigidity,def:klein_bound,thm:recon_monodromy_correct}
  A group exceeding the Klein cap is forced off genus 0; the threshold gap is
  then bounded by $2\genus$. The dichotomy: either genus-0 with a bounded group
  and no gap, or positive genus with a gap.
\end{theorem}

\part{The instances}

\chapter{den Boer: the perfect five-card MPC}

\begin{definition}[The five-card scheme and plug]
  \label{def:five_card}
  \rocq{pgg_smc.five_card_scheme_I5.fcI_scheme,
       pgg_smc.five_card_scheme_I5.fcI_reconK,
       pgg_smc.five_card_kim.FiveCardKim_M,
       pgg_smc.five_card_family.five_card_plug,
       pgg_smc.five_card_family.fcI_perm_compatible_kim}
  \rocqok
  \uses{def:threshold_scheme,def:recon_plug,def:monodromy_repr}
  The five-card scheme reconstructs the secret bit as ``three consecutive
  hearts''. Its monodromy is the cyclic group $C_5$ generated by the 5-cycle,
  and reconstruction is invariant over the whole group; together they form
  $\mathtt{five\_card\_plug}$.
\end{definition}

\begin{definition}[The den Boer encoding]
  \label{def:den_boer_encoding}
  \rocq{pgg_smc.den_boer_encoding.den_boer_layout,
       pgg_smc.den_boer_encoding.den_boer_encoding,
       pgg_smc.den_boer_encoding.den_boer_orbit_perm}
  \rocqok
  \uses{def:input_encoding,def:five_card}
  The two committed bits $a,b$ determine the starting layout, which is a valid
  sharing of $a\AND b$, and equal-output inputs are one cyclic rotation apart.
  This instantiates the input encoding at $f(a,b)=a\AND b$.
\end{definition}

\begin{theorem}[The running protocol computes $a\AND b$]
  \label{thm:den_boer_run_output}
  \rocq{pgg_smc.den_boer_run.den_boer_run_output,
       pgg_smc.den_boer_run.den_boer_dealer_layout,
       pgg_smc.den_boer_run.den_boer_decodeK}
  \rocqok
  \uses{def:den_boer_encoding,thm:recon_monodromy_correct}
  The committed dealer injects the input-derived layout through the content
  readout, and recovery returns $a\AND b$ for every cut. Zero axioms.
\end{theorem}

\begin{lemma}[The layout dealer is dual]
  \label{lem:den_boer_duality}
  \rocq{pgg_smc.den_boer_run.den_boer_layout_player0_dual,
       pgg_smc.den_boer_run.den_boer_layout_verifier_dual}
  \rocqok
  \uses{thm:den_boer_run_output,lem:session_duality}
  The input-derived-content dealer keeps the session type of the constant
  dealer, so all four duality checks discharge by reduction.
\end{lemma}

\begin{theorem}[Perfect input privacy]
  \label{thm:den_boer_input_private}
  \rocq{pgg_smc.den_boer_encoding.den_boer_input_private,
       pgg_smc.den_boer_encoding.den_boer_cinde}
  \rocqok
  \uses{def:den_boer_encoding}
  The conditional mutual information $\MI(\text{inputs};\text{view}\mid
  \text{output}) = 0$: a coalition learns nothing about the inputs beyond
  $a\AND b$, because equal-output inputs are cyclic rotations giving an
  identically distributed view under the uniform cut.
\end{theorem}

\begin{theorem}[The partial-view leakage ramp]
  \label{thm:leak_ramp}
  \rocq{pgg_smc.five_card_leakage.leak_k1,
       pgg_smc.five_card_leakage.leak_k2_adj,
       pgg_smc.five_card_leakage.leak_k2_dist2,
       pgg_smc.five_card_leakage.leak_k3,
       pgg_smc.five_card_leakage.leak_k4,
       pgg_smc.five_card_leakage.leak_k5}
  \rocqok
  \uses{def:den_boer_encoding}
  A coalition seeing $k$ cards leaks a quantified amount of the secret: nothing
  at $k=1$, a closed-form positive amount at $k=2$ that depends on whether the
  cards are adjacent or distance-2, growing to the full $H(\tfrac14)$ at $k\ge4$.
\end{theorem}

\begin{theorem}[Profile and end-to-end correctness]
  \label{thm:den_boer_profile}
  \rocq{pgg_smc.den_boer_profile.den_boer_profile,
       pgg_smc.den_boer_profile.den_boer_perfect,
       pgg_smc.den_boer_profile.FiveCardKim_protocol_correct,
       pgg_smc.den_boer_profile.profile_k_denboer}
  \rocqok
  \uses{def:monodromy_profile,def:five_card,thm:recon_monodromy_correct}
  den Boer is the $\eps=0$ member of the five-card family: the dealing phase is
  perfectly anonymous, the privacy threshold is two, and reconstruction recovers
  the secret bit for any cut.
\end{theorem}

\chapter{The five-card family and biased anonymity}

\begin{definition}[The five-card family]
  \label{def:five_card_family}
  \rocq{pgg_smc.five_card_family.five_card_profile,
       pgg_smc.rigidity_kim_instance.kim_profile}
  \rocqok
  \uses{def:monodromy_profile,def:five_card}
  The family is parameterized by the cut bias $\eps$; den Boer is the $\eps=0$
  (perfect) member, and Kim's instance is the biased one.
\end{definition}

\begin{theorem}[Spectral-gap anonymity]
  \label{thm:kim_mixing}
  \rocq{pgg_smc.five_card_kim.kim_lambda2,
       pgg_smc.five_card_kim.kim_spectral_gap,
       pgg_smc.five_card_kim.kim_spectral_convergence,
       pgg_smc.five_card_kim.kim_var_dist_exact,
       pgg_smc.five_card_kim.fc_kim_security_witness}
  \rocqok
  \uses{def:five_card_family,def:security_witness}
  Kim's Schreier walk is doubly stochastic with second eigenvalue
  $\lambda_2=\tfrac54|\eps|$; the spectral gap $1-\lambda_2$ is positive and
  drives an exact variation-distance bound. This exercises the security-witness
  interface with a genuine mixing argument, of which den Boer's $\eps=0$ case is
  the sharp limit.
\end{theorem}

\chapter{The $S_5$ no-go: a representation-theoretic obstruction}

The admissibility gate of the previous chapter rejects a monodromy group by a
counting argument: a group too large for the Klein cap cannot act on a genus-0
curve, so its threshold gap is forced wide. That argument bounds how large the
gap must grow. It does not, on its own, show that a usable scheme fails to exist.
This chapter proves the sharper statement for the wired $S_5$ instance. It
exhibits a single representation-theoretic obstruction, internal to the group and
independent of any curve, that forbids the secret-encoding code outright. The
whole no-go collapses to one fact about the submodule lattice of the natural
permutation module of $S_5$ in characteristic five.

We first fix how the secret and the shares sit inside a vector space.

\begin{definition}[The wired $S_5$ representation]
  \label{def:wired_s5}
  \rocq{pgg_reconstruct.s5_nogo.rG_secret,
       pgg_reconstruct.s5_nogo.e0,
       pgg_reconstruct.s5_nogo.proj_share}
  \rocqok
  The wired instance places the secret in the fixed coordinate $0$ of
  $\GF(5)^6$ and the five shares in the coordinates $1$ through $5$, permuted by
  $S_5$. The secret direction is the basis vector $e_0$ at coordinate $0$.
\end{definition}

A covering scheme recovers the secret by reading it off a linear code that the
monodromy leaves invariant. Concretely, a recovering code is an $S_5$-submodule
$U$ of $\GF(5)^6$ subject to two demands. It must be invariant, because the cut
acts through the monodromy and the recovered value may not depend on which cut
was drawn. And it must contain the secret direction $e_0$, because the secret is
the value stored in coordinate $0$. The existence of a recovering code is thus
exactly the existence of a secret-carrying invariant submodule.

Not every such submodule yields a threshold scheme. If $U$ is all of
$\GF(5)^6$, every coordinate is needed to pin down a codeword and there is no
threshold at all. If $U$ is only the line through $e_0$, no genuine sharing takes
place. A non-degenerate threshold gap lives strictly between these extremes, and
the algebraic-geometry dimension count makes the admissible band precise. At the
six evaluation points of the wired instance, the Riemann-Roch relations of an
AG-Massey code force a strict gap to come from a code whose dimension is neither
too small nor too large.

\begin{definition}[The gap-dimension window]
  \label{def:gap_window}
  \rocq{pgg_reconstruct.invariant_profiler.secret_inv_dim,
       pgg_reconstruct.invariant_profiler.feasible,
       pgg_reconstruct.gap_dimension.gap_dim_window}
  \rocqok
  \uses{def:covering_scheme}
  Under the AG-Massey relations at six evaluation points, a strict threshold gap
  forces the secret-encoding code dimension into the window $\{3,4\}$.
\end{definition}

The secret coordinate plays a distinguished role: it is the one coordinate the
group fixes. Forgetting it should therefore lose nothing about the group action
and exactly one dimension about the secret, which is what makes the
six-dimensional problem collapse to a five-dimensional one. Write $\mathrm{Pr}$
for the projection that drops coordinate $0$ and keeps the five share
coordinates. It intertwines the wired action with the natural permutation action,
so it carries an invariant submodule to an invariant submodule of the natural
module $\GF(5)^5$. Its kernel is precisely the secret line, so a submodule that
contains $e_0$ meets that kernel in one dimension and loses exactly that one
dimension under projection.

\begin{lemma}[The reduction: projection drops the dimension by one]
  \label{lem:reduction}
  \rocq{pgg_reconstruct.s5_nogo.proj_mxmodule,
       pgg_reconstruct.s5_nogo.mxrank_proj_pred,
       pgg_reconstruct.s5_nogo.e0_proj_share,
       pgg_reconstruct.s5_nogo.proj_share_rank}
  \rocqok
  \uses{def:wired_s5}
  The projection $\mathrm{Pr}$ sends a secret-carrying invariant submodule of
  $\GF(5)^6$ of dimension $d$ to an invariant submodule of the natural module
  $\GF(5)^5$ of dimension $d-1$.
\end{lemma}

Apply the reduction to the gap window. A secret-encoding code of dimension $3$
projects to an invariant submodule of the natural module of dimension $2$, and a
code of dimension $4$ projects to one of dimension $3$. The whole question is now
internal to the natural permutation module of $S_5$: does it carry an invariant
subspace of dimension $2$ or $3$? The answer is no, for a reason special to
characteristic five.

\begin{theorem}[The kernel fact: no submodule of dimension $2$ or $3$]
  \label{thm:kernel_fact}
  \rocq{pgg_reconstruct.s5_nogo.perm_module_no_dim23}
  \rocqok
  \uses{def:wired_s5}
  Every $S_5$-submodule of the natural module $\GF(5)^5$ has dimension in
  $\{0,1,4,5\}$; none has dimension $2$ or $3$.
\end{theorem}

The proof is elementary and turns on the difference vectors $e_i - e_j$. The
natural module has two evident invariant subspaces: the line spanned by the
all-ones vector and the hyperplane of sum-zero vectors. The claim is that nothing
lies strictly between dimension one and dimension four. Let $W$ be a submodule.
Either every vector of $W$ is constant, in which case $W$ lies in the all-ones
line and has dimension at most one, or $W$ contains a vector $v$ with two
distinct coordinates $v_i \ne v_j$. In the second case, subtract from $v$ its
image under the transposition that swaps $i$ and $j$. The difference is a nonzero
scalar multiple of $e_i - e_j$, so after rescaling $e_i - e_j$ itself lies in
$W$. The group $S_5$ is doubly transitive, so it carries this one difference
vector to every $e_a - e_b$, and those vectors span the whole sum-zero subspace,
of dimension four. Hence $W$ has dimension at least four. The middle dimensions
two and three never occur.

Characteristic five is not a convenience here, it is the mechanism. When the
field characteristic is coprime to the group order, Maschke's theorem splits
every module into irreducible summands, and the sum-zero subspace breaks as the
all-ones line plus a three-dimensional irreducible heart. That heart would be a
submodule of dimension three, and the forbidden dimensions would reappear. The
obstruction survives precisely because five divides the dimension five, which
forces the all-ones vector into the sum-zero subspace and collapses the lattice
into a single chain.

\begin{theorem}[A characteristic-five contrast]
  \label{thm:maschke_contrast}
  \rocq{pgg_reconstruct.invariant_profiler.maschke_ss}
  \rocqok
  \uses{thm:kernel_fact}
  When the characteristic is coprime to $|G|$, Maschke's theorem makes the module
  completely reducible, so the heart of dimension three appears and the kernel
  fact fails. The no-go is a modular phenomenon.
\end{theorem}

Everything now assembles. A gap instance needs a secret-encoding code of
dimension three or four. By the reduction such a code yields an invariant
submodule of the natural module of dimension two or three. The kernel fact
forbids both. There is therefore no secret-encoding invariant code in the gap
window, and the wired $S_5$ gap is mathematically impossible.

\begin{theorem}[The $S_5$ no-go]
  \label{thm:s5_nogo}
  \rocq{pgg_reconstruct.s5_nogo.s5_no_secret_dim3,
       pgg_reconstruct.s5_nogo.s5_no_secret_dim4,
       pgg_reconstruct.s5_nogo.s5_gap_window_infeasible,
       pgg_reconstruct.s5_nogo.s5_gap_infeasible}
  \rocqok
  \uses{thm:kernel_fact,lem:reduction,def:gap_window}
  No secret-encoding $S_5$-invariant code of dimension $3$ or $4$ exists, so the
  wired $S_5$ gap is infeasible: $S_5$ admits no AG code in the gap window. The
  result is axiom-free.
\end{theorem}

Read through the covering dictionary, an invariant code of the right dimension is
an algebraic-geometry code on the covering curve, and the no-go says $S_5$
supports none in the gap window. The only realizations available are forced up to
positive genus, where the gap exceeds the number of shares and the threshold
becomes vacuous. This is one rung of a ladder that the framework makes explicit:
genus zero gives the exact Shamir threshold with no gap, and each unit of genus
widens the gap by two. A group as large as $S_5$ is pushed past the rung where a
threshold scheme is still useful.

\begin{definition}[The genus ladder]
  \label{def:genus_ladder}
  \rocq{pgg_reconstruct.cover_genus0.shamir_exact,
       pgg_reconstruct.cover_genus1.higher_genus_gap_bound,
       pgg_reconstruct.cover_genus2.genus2_gap,
       pgg_reconstruct.combinatorial_rigidity.CombinatorialRigidity}
  \rocqok
  \uses{def:covering_scheme}
  Genus $0$ gives an exact (Shamir) threshold; genus $g$ gives a gap of $2g$
  (elliptic $2$, hyperelliptic $4$). Combinatorial rigidity is the curve-free
  dual, replacing the genus-0 cap with the order inequality $|G| >
  \mathtt{klein\_genus0\_bound}$.
\end{definition}

\chapter{$S_5\times S_5$ and the Klein-cap genus context}

\begin{definition}[The $S_5$ instance]
  \label{def:s5_instance}
  \rocq{pgg_smc.rigidity_s5_instance.s5_rigidity,
       pgg_smc.rigidity_s5_instance.s5_brings_covering_genus,
       pgg_smc.rigidity_s5_instance.s5_genus0_klein,
       pgg_smc.rigidity_s5_instance.s5_tradeoff}
  \rocqok
  \uses{def:algebraic_rigidity,def:klein_bound,def:genus_ladder}
  $|S_5|=120$ exceeds the Klein cap, so genus 0 is impossible; the instance
  lands on Bring's curve at genus $4$, with a gap that exceeds $N$ and is
  therefore vacuous. The rigidity witness exists only at this cost.
\end{definition}

\begin{definition}[$S_5$ group order]
  \label{ax:s5_group_order}
  \rocq{pgg_smc.rigidity_s5_instance.s5_group_order_eq}
  \uses{def:s5_instance}
  Assumed: $|G| = 120$ for the path-generator $S_5$. The kernel computation of
  the generated subgroup order is infeasible, so this is a justified axiom.
\end{definition}

\begin{definition}[Bring's curve realization]
  \label{ax:s5_brings}
  \rocq{pgg_smc.rigidity_s5_instance.s5_brings_covering_realised,
       pgg_smc.s5_mixing.s5_rayleigh_Q2_R}
  \uses{def:s5_instance}
  Assumed: Bring's curve realizes the genus-4 covering data (Edge 1978), and a
  Rayleigh bound on the squared transition operator (an externally certified
  sum-of-squares certificate). Cited classical / external facts.
\end{definition}

\begin{definition}[The $S_5\times S_5$ instance]
  \label{def:s5x5_instance}
  \rocq{pgg_smc.rigidity_s5x5_instance.s5x5_rigidity,
       pgg_smc.rigidity_s5x5_instance.s5x5_protocol_correct,
       pgg_smc.rigidity_s5x5_instance.s5x5_combinatorial_rigidity,
       pgg_smc.rigidity_s5x5_instance.s5x5_tradeoff,
       pgg_smc.s5x5_pile.s5x5_pile1_stab}
  \rocqok
  \uses{def:algebraic_rigidity,def:genus_ladder}
  The product $S_5\times S_5$ ($|G|=14400$) also exceeds the cap and is forced
  to positive genus; the protocol's reconstruction is still correct, but the
  scheme is a combinatorial-rigidity (curve-free) instance, not a sharp AG code.
\end{definition}

\begin{definition}[$S_5\times S_5$ realizations]
  \label{ax:s5x5}
  \rocq{pgg_smc.rigidity_s5x5_instance.s5x5_group_order_eq,
       pgg_smc.rigidity_s5x5_instance.s5x5_inverse_galois_realised,
       pgg_smc.rigidity_s5x5_instance.s5x5_multi_realised}
  \uses{def:s5x5_instance}
  Assumed: $|G|=14400$; a genus-173 Galois cover of $\mathbb{P}^1$ with deck
  group $S_5\times S_5$ exists (inverse Galois); and the two-component Bring
  realization exists. Cited classical facts.
\end{definition}

\chapter{The modularity ledger}

Each interface slot of Part~I is filled by exactly one node per instance, and
the generic theorems are inherited rather than reproved. den Boer fills all six
slots (Figure~\ref{fig:bp:modularity}) and is the only instance carried to a
full MPC end-to-end. Kim fills the monodromy and security-witness slots with a
biased, spectral fill. $S_5$ and $S_5\times S_5$ reach the admissibility gate
and are rejected as AG codes by the Klein cap, the first through the
representation-theoretic no-go (Theorem~\ref{thm:s5_nogo}), the second through
the order inequality of combinatorial rigidity. The generic side is entirely
axiom-free; the six assumptions live only in the geometric realization of the
two large groups.
